% Generated draft: requires mathematical and source review.
\documentclass[11pt]{article}
\usepackage[margin=25mm]{geometry}
\usepackage{fontspec}
\setmainfont{Latin Modern Roman}
\usepackage{amsmath,amssymb,mathtools,unicode-math}
\setmathfont{Latin Modern Math}
\usepackage{enumitem,needspace,xurl,hyperref}
\setlength{\parindent}{0pt}
\setlength{\parskip}{5pt}
\setlength{\emergencystretch}{3em}
\newcommand{\N}{\mathbb{N}}
\newcommand{\Z}{\mathbb{Z}}
\newcommand{\Q}{\mathbb{Q}}
\newcommand{\R}{\mathbb{R}}
\newcommand{\C}{\mathbb{C}}
\newcommand{\F}{\mathbb{F}}
\newcommand{\A}{\mathbb{A}}
\newcommand{\PP}{\mathbb{P}}
\newcommand{\E}{\mathbb{E}}
\newcommand{\eps}{\varepsilon}
\newcommand{\dd}{\,\mathrm{d}}
\begin{document}
\section*{UnsolvedMath ID: 30005754}
\textbf{Source problem code:} catalogue entry
\subsection{Definitions and mathematical statement}
The supplied source denotes a cluster algebra by \(k[A]\) and refers to a dual canonical basis \(B^*\) in a reference Lie-theoretic setting. It also uses an object \(N\) and an inclusion condition \(A\subseteq N\). The record does not specify the coefficient field or ring \(k\), the objects \(A\) and \(N\), the ambient space in which this inclusion is interpreted, or the algebra having basis \(B^*\). These symbols therefore retain their source-dependent meanings; no additional hypotheses are imposed here.

\textbf{Research problem.} For general cluster algebras denoted by \(k[A]\), identify appropriate analogues of the dual canonical basis \(B^*\). Formulate a corresponding generalization of the Fomin--Zelevinsky conjecture that includes the source's intended cases
\[
A\not\subseteq N.
\]

\textbf{Unresolved formulation qualifications.} The record does not state the Fomin--Zelevinsky conjecture being generalized or specify which properties characterize an acceptable analogue of \(B^*\). It also does not specify the class of cluster algebras, coefficient conventions, or whether the desired bases belong to classical or quantum algebras. Consequently, the supplied evidence determines a research question but does not determine a precise quantified conjecture or a testable list of basis axioms. In particular, the condition \(A\not\subseteq N\) cannot be interpreted independently of the missing ambient setting.
\subsection{Short English statement}
Find analogues of the dual canonical basis for general cluster algebras and extend the intended Fomin--Zelevinsky conjecture beyond the setting \(A\subseteq N\). The supplied record leaves the underlying objects, basis requirements, and conjecture unspecified, so a fully precise formulation requires further source context.
\subsection{Sources}
\begin{itemize}
\item Supplied UnsolvedMath local registry record: permanent ID \(30005754\), title \emph{Dual Canonical Basis Analogues for General Cluster Algebras}, exact source problem code \texttt{catalogue entry}; collection: Oberwolfach Reports Open Problems; proposed year: 2024.
\item Supplied report citation: \emph{Cluster Algebras and Its Applications} (2024), Oberwolfach Reports, pp.\ 69--136, DOI 10.4171/owr/2024/2. Original OWR report: \url{https://doi.org/10.4171/owr/2024/2}. The supplied background describes these reports as documenting mathematical workshops and their open-problem discussions.
\item The original extracted statement associates the question with \emph{Analogs of dual canonical bases for cluster algebras from Lie theory} and Fan Qin. The separate proposer field is empty, so this association is retained without asserting that Fan Qin formally proposed the problem.
\item Supplied literature source, labelled ``Lie-theoretic analogs paper'': \url{https://arxiv.org/abs/2407.02480}.
\item Supplied literature source, labelled ``quantum-cluster basis paper'': \url{https://arxiv.org/abs/2003.13674}.
\item Literature assessment date: \texttt{2026-08-22}. This is supplied dated curation derived from the OWR curation, not independently verified here. The record labels the problem partially solved and assesses that Qin constructs dual-canonical or triangular analogues in broad Lie-theoretic families, while analogues for arbitrary cluster algebras remain unresolved. The source explicitly requests independent verification of time-sensitive status and cited claims.
\item The record labels its statement \texttt{corrected\_verified} and reports that the official report was checked directly and that the cleaned statement restores definitions and removes extraction artifacts or source-dependent wording. This is a supplied source-verification claim, not an independent verification in this draft; the supplied text still omits the definitions and conjecture needed for a precise formulation.
\end{itemize}
\end{document}
